\chapter{Expanded measurement, verification and coefficient evidence}
\label{ch:expand-verification}
\section{Verification, validation, selection and qualification are different}
Verification asks whether an implementation performs the stated calculation.
A finite-difference derivative, an integral-image hand sum and comparison with
a separately implemented operation are examples. Validation asks whether the
method represents the intended task: correct line ink, readable letters or
actual carriage distances. Model selection chooses settings using observed
performance; qualification assesses the specified requirement under the
documented test conditions. A high validation score used to choose a width
is evidence for selection, not an untouched qualification result.

The additions explain how calculations can be checked; none was rerun as an
experiment. New arithmetic is labelled as constructed or derived. No coefficient
sweep, inference, robot motion, training, lighting experiment, old-segmenter
comparison, pressure repair or demonstration preparation belongs to this work.

\section{How the recorded NumPy--Torch comparisons should be understood}
The final report Table 7 records reader/writer maximum-logit differences
normalised by maximum magnitude of $2.12\times10^{-14}$ and
$3.04\times10^{-15}$; absolute loss differences of $1.14\times10^{-13}$
and $1.78\times10^{-15}$; worst-ranked gradient relative L2 differences of
$2.80\times10^{-13}$ and $6.01\times10^{-14}$~\cite{finalreport2026}.
Its prose saying all values are below $10^{-13}$ is inconsistent with two
displayed reader entries. The table values are retained rather than replaced
with that blanket description.

For corresponding tensors $a,b$, useful measures are
\begin{align*}
e_{\max}&=\max_i|a_i-b_i|,\\
e_{\rm rel,max}&=e_{\max}/\max(\epsilon,\max_i|b_i|),\\
e_{\rm rel,2}&=\|a-b\|_2/\max(\epsilon,\|b\|_2).
\end{align*}
Other symmetric denominators are possible and give different numbers.
The exact archived checker defines which was used; a descriptive table
heading alone does not recover every denominator/floor. Relative metrics
near zero can be misleading, so retain an absolute metric too. ``Worst ranked''
must identify the compared parameter tensors, not imply all possible states
and inputs were exhausted.

A meaningful comparison holds fixed input arrays, parameters, precision,
layout, padding, activation conventions, training/evaluation mode, batchnorm
statistics, LSTM gate order, blank index, sequence lengths and loss reduction.
Dropout is disabled or its mask is shared. Compare intermediate operations
before relying on final logits. Float64 agreement near these scales supports
consistency for the checked cases; it does not validate the model's learned
accuracy or guarantee a whole training trajectory matches. Floating-point
operations are order-dependent~\cite{torchnumeric}. Independently trained
checkpoints can differ for many reasons beyond accumulated rounding error.

Table 6's performance comparison is a different experiment: reader pooled
94.06\% versus 94.39\%, writer 98.08\% versus 97.51\% on common material.
These are trained-route evaluation scores, not numerical roundoff measures.
Matching an operation to Torch and achieving a character score answer
different questions. The writer QTP2 score 153/156 should not be silently
assigned to the alternative first-principles checkpoint.

\section{Finite differences and all Appendix A.1 arithmetic}
For scalar parameter $\theta_i$, central difference is
\[
g_i^{\rm FD}={L(\theta+h e_i)-L(\theta-h e_i)\over2h}.
\]
Taylor expansion cancels the first-order error terms and gives an $O(h^2)$
truncation error when $L$ is sufficiently smooth. Making $h$ too small
causes subtraction cancellation and roundoff; making it too large measures
curvature rather than the local derivative. ReLU zeros and max-pool ties
are nondifferentiable, so a central difference crossing them need not equal
the implementation's selected subgradient. These observations follow from
calculus and floating-point arithmetic, not new network tests.

The report records losses 0.4703364611173033 and 0.4703403997759956,
$h=10^{-4}$. In the recorded plus/minus order,
\[
{0.4703364611173033-0.4703403997759956\over0.0002}
\simeq-0.0196932934615.
\]
The analytic derivative is $-0.0196932934309$, absolute discrepancy
about $3.06\times10^{-11}$. This is one recorded coordinate check,
not a worst-case bound for all derivatives. Recovering it operationally
would require the exact frozen input, target, parameter coordinate, dtype,
weights, loss reduction and perturbation order; the two scalar losses alone
do not identify those objects. No checker was executed in this edition.

The patch-storage example is $368640\times4/2^{20}=1.40625$ MiB.
The factor four assumes four bytes per element (float32); float64 would
double the value. MiB uses $2^{20}$ bytes, while MB uses $10^6$.
An im2col view can become a materialised array during reshape/multiplication;
count actual storage, not only logical view shape. This arithmetic does not
establish peak process memory. The switch-confirmation example is
$(3-1)\times6=12$ ms; three reads contain two intervals.

\section{The mysterious ``Appendix calculation 1'': the page vote}
The final report labels Appendix A.0.1 ``System integration calculation''.
It combines line writer probabilities using ink-count weights:
\[
\bar p_k={\sum_i w_i p_{i,k}\over\sum_i w_i},\quad w_i\ge0.
\]
For two archived Author O01 crops,
\[
\bar p_O={20204(0.9996674)+21750(0.9999279)\over41954}
\simeq0.9998024.
\]
Weights sum to 41954; normalised weights sum to one, so the pooled vector
remains a probability vector when the inputs are. Weighting by ink count
assumes more ink supplies more useful evidence; a dark artefact or a long
misclassified line can instead dominate. This is not a Bayesian product
of independent likelihoods, and near-one softmax values do not prove
near-certain real-world identity. The example uses two crops and cannot
substitute for the full-page result~\cite{finalreport2026}.

\section{Levenshtein edits, character accuracy and exact text measures}
For reference prefix length $i$ and predicted prefix length $j$,
\[
D_{i,j}=\min\{D_{i-1,j}+1,D_{i,j-1}+1,
D_{i-1,j-1}+\mathbf1[t_i\ne\hat t_j]\},
\]
with $D_{i,0}=i,D_{0,j}=j$. The three branches represent deletion,
insertion and substitution/match. This dynamic-programming calculation gives
minimum edit count~\cite{levenshtein1966}. For ``cat'' versus ``cut'',
one substitution gives CER $1/3$ and character score $2/3$. A missing
space counts if spaces are included. Case folding and punctuation removal
must be specified before comparison; changing them changes the task.

Pooled CER is $\sum_iD_i/\sum_i|t_i|$, not the unweighted mean of line
CERs. For lengths 2 and 98 with errors 1 and 0, pooled accuracy is 99\%,
whereas mean line accuracy is 75\%. The final report uses the pooled
formula. Its raw accuracy can become negative if insertions exceed the
reference length; the candidate \texttt{CharAcc} clamps at zero. These are
different conventions at that extreme.

Exact-line match counts a line only if the entire chosen text agrees.
Word error rate, fraction of exact words and longest-common-subsequence
word recall are also different. The inspected \texttt{WordAcc} computes
case-insensitive word-level LCS recall, despite its short descriptive comment.
It should not automatically be interpreted as standard WER complement or
strict positional exact-word accuracy. Reported numerical evidence must be
read with the actual scorer and text normalisation.

\section{Wilson intervals: derivation, numerical meaning and dependence}
For $x$ correct binary decisions out of $n$ and $\hat p=x/n$, invert the
score inequality $(\hat p-p)^2\le z^2p(1-p)/n$. Solving the resulting
quadratic gives
\[
{\hat p+z^2/(2n)\ \pm\ z\sqrt{\hat p(1-\hat p)/n+z^2/(4n^2)}
\over1+z^2/n}.
\]
Use $z\simeq1.96$ for a nominal two-sided 95\% interval~\cite{wilson1927,nistinterval}.
For $153/156$, $\hat p\simeq0.98077$ and the interval is about
94.5--99.3\%; for $159/184$, about 80.7--90.6\%. These are reported
line-decision examples. Even 25/25 gives a lower bound about 86.7\%, not
certainty of universal success.

Wilson avoids the zero-width interval that a naive normal approximation
can give at all successes. Its binomial interpretation still assumes suitable
independent trials with a shared success probability. Lines sharing page,
writer, lighting and model are dependent. The report therefore treats its
intervals descriptively. It is incorrect to claim a 95\% posterior probability
that a fixed unknown $p$ lies in a frequentist interval without an additional
Bayesian model. Character edit scores are not simple independent Bernoulli
trials and cannot be inserted as fractional ``successes'' without justification.

\section{Direct geometry: registration before comparison}
The direct word-shape check is independent of a reader's text labels at
scoring time, but depends on correct crop/transcript matching and image
processing. Binary masks are thinned to centreline point sets. Normalisation
and registration define which differences are ignored. Translation removes
position; isotropic scaling removes uniform size; independent X/Y scaling
also removes aspect ratio; shear correction removes slant. A deformable warp
can erase precisely the writer features being assessed. Therefore allowable
transforms are part of the metric, not cosmetic preparation.

The archived natural-baseline protocol documents two alternatives: line
x-height scaled to 48 pixels, or tightly cropped whole-word ink height scaled
isotropically to 48 pixels; centring and at most five pixels translation per
axis; point tolerance 4.8 pixels. Whole-word height includes ascenders and
descenders, unlike body x-height. No width-only stretch, slant correction or
deformable warp belongs to that protocol. A 4.8-pixel tolerance is one tenth
of the selected 48-pixel reference height; changing that reference changes
the physical meaning of tolerance. The local evidence source is
\texttt{Report/NaturalVariationBaseline\_2026-11-02/PROTOCOL.md} and
\texttt{FINDINGS.md}. Its historical aggregate is not substituted for the
final report's later aggregation.

Optimising translation for maximum overlap already introduces an alignment
selection step. The same policy must be applied to genuine and robot pairs.
Reference crops must be exact-token matches and self-matches excluded.
Nearest-reference selection gives a higher or equal score than averaging
every reference and becomes easier with more references. Matching reference
banks controls that opportunity. A plot of beautifully aligned skeletons
must retain the raw word photographs and transformation policy for context.

\section{Point-set precision, recall and harmonic F-score}
Let $A$ be robot skeleton points and $B$ genuine-reference points. With
distance tolerance $\tau$, define
\[
P={|\{a\in A:\min_{b\in B}\|a-b\|\le\tau\}|\over|A|},\qquad
R={|\{b\in B:\min_{a\in A}\|b-a\|\le\tau\}|\over|B|},\quad
F={2PR\over P+R}.
\]
The thresholded point-set definition is used by reconstruction benchmarks;
its application to 2D handwriting and the chosen normalisation are project
choices~\cite{tanksmetric}. Each point asks whether any nearby point exists;
this is not one-to-one matching. A distance transform gives distance to the
nearest foreground grid point; a spatial tree can independently compute
nearest-point distances for the same sets. Agreement checks arithmetic,
not semantic validity of the chosen crops.

Constructed example: 90 of 100 robot points have a reference neighbour,
while 80 of 100 reference points have a robot neighbour. Then
$P=0.9,R=0.8,F=1.44/1.7\simeq0.8471$. Missing strokes lower recall;
extra strokes lower precision. F is not intersection-over-union, character
correctness or Pearson correlation. Tolerance can make shifted strokes
match despite zero exact pixel overlap. If both sets are empty, or one is
empty, a scoring convention is required; silent division by zero is not a
definition. Scores can depend on point density unless sampling is controlled.

\section{Natural-variation normalisation and why 100\% is not a ceiling}
For writer $k$, let $F_{r,k}$ and $F_{g,k}$ be robot and genuine agreement
against corresponding reference banks, using the same protocol. The final
report defines $Q_k=100F_{r,k}/F_{g,k}$ and averages eleven writer ratios.
For writer 150 the displayed $87.73/90.01$ gives about 97.47\%.
For Owen $70.22/76.03$ gives about 92.36\%, subject to rounded table inputs.
The final result is about 86.7\% relative agreement, not 86.7\% correctly
written characters~\cite{finalreport2026}.

Equal-writer mean $\frac1K\sum_kF_{r,k}/F_{g,k}$ differs from
ratio of means $\sum_kF_{r,k}/\sum_kF_{g,k}$. For two writers with
$(F_g,F_r)=(0.9,0.9),(0.5,0.25)$, mean ratio is 0.75 but ratio of
means is $1.15/1.4\simeq0.8214$. This explains why an older evidence
note can quote a different aggregate without either arithmetic necessarily
being wrong. State the exact rule rather than mixing versions.

If $F_{r,k}>F_{g,k}$, $Q_k$ exceeds 100\%; it means closer agreement to
the available bank than the measured genuine baseline, not superhuman style.
A small denominator destabilises the ratio. A baseline below one does not
make an absolute 85\% target impossible; the ratio is a newly defined
comparison, not a logical replacement for the proposal's original requirement.
Repeated words are a selected subset, often shorter/common words. Four or
five words for a writer do not establish full-page fidelity. A nearest-reference
figure can have a higher F than the corresponding ordinary-reference writer
mean. Both policies should be named.

\section{Direct length checks and completion-time statistics}
For commanded length $L_0$ and measured $L_i$, signed error is $e_i=L_i-L_0$;
absolute error $|e_i|$; mean signed error indicates bias; mean absolute error
describes typical magnitude; maximum absolute error addresses a worst
observed limit. The report's 70 readings give mean absolute error 0.297 mm
and maximum 0.49 mm. Neither is an uncertainty bound for all future drawings.
A ruler/caliper's resolution, endpoint definition, paper motion and repeated
measurements affect interpretation. A sloped line length may combine errors
in both axes and squareness. Scaling calibration and independent acceptance
measurements should not be confused.

For completed job $j$, normalised completion time is $r_j=t_j/n_j$, where
the counting convention for $n_j$ must include/exclude spaces consistently.
The mean of job ratios $\frac1J\sum_jr_j$ gives every job equal weight;
$\sum_jt_j/\sum_jn_j$ weights by characters. For 10 characters in 30 s
and 100 in 100 s, these give 2 s/character and about 1.182 s/character.
The report's 25 confirmed completed jobs and recorded stages must be kept
distinct from unconfirmed/pending rows and ideal trajectory budgets.
Full request time includes generation and motion when the record says so;
classifier-only latency does not. Host names and laptop/ODROID evidence
must remain attached to the timings.

The conditional 2.76 network-speed factor from a separate benchmark is not
a measured system acceleration. If image time is $t_I$ and selected network
time $t_N$, the conditional estimate is $t_I+t_N/2.76$ plus all unchanged
stages. Dividing the entire request by 2.76 would overstate the gain. This
is Amdahl-style accounting of unchanged work, not a new passing qualification.

\section{Capture records, plots and the old-segmenter question}
The report's Table 18 lists expected/reduced/gradient/+3-degree/90-degree
line counts: Owen $20/19/20/11/19$, Dylan $24/24/24/25/24$, and
Yeukita $26/26/26/26/26$. These are digital perturbation observations,
not an illumination robustness guarantee. Equal counts do not establish
correct labels, ownership or readable crops. Reduced resolution can merge
thin strokes; a synthetic gradient is not every physical shadow, glare or
phone-light configuration.

Seam-sensitivity plots vary a cost coefficient and show line/ownership
behaviour. A visually cleaner seam is evidence for that example, not a
universal coefficient optimum. CNN channel plots show activations under one
checkpoint/input; independent colour scales prevent comparing absolute
strength across panels. CTC plots show frame symbols and constrained paths;
Bresenham plots show integer command locations. Joint-score plots show
model-based selection, not independent physical style accuracy.

The requested old-versus-new segmentation accuracy comparison is a meaningful
paired-ablation question, but running it is explicitly outside this task.
To interpret a future result, the same frozen pages, labels, classifier,
preparation and scoring would be needed. Failed or unmatched crops must
remain accounted for; evaluating only surviving easy crops can inflate the
score. A report saying segmentation ``affects accuracy'' is plausible causal
reasoning, but a quantified benefit requires that controlled evidence.

\section{Coefficient evidence register and questions left open}
\begin{description}[style=nextline]
\item[Derived from stated geometry/electronics.] Forty mm/rev, nominal
80 steps/mm, 12.5 micrometre spacing, nearest-rounding bounds, RPM conversion,
12 ms confirmation span, conditional reference voltage and resistor currents.
These validate arithmetic only under their stated inputs.
\item[Published method, project constants.] Otsu and integral sums;
deskew objective; seam dynamic programming; Zhang--Suen; CTC and LSTM;
dropout; Douglas--Peucker; optimiser families; edit distance and score interval.
Their publications do not establish the project's exact parameter choices.
\item[Reported selection evidence.] Stroke-width 0.17 versus the other
three displayed widths. Keep the unresolved page/line wording and selected
validation status visible. Do not promote it to an independent result.
\item[Documented local rationale, not newly verified.] 0.45 connection
suppression between commented print/joined examples; 15\%/threefold solid-crop
guard with commented crop-density observations; visual settling dwell.
These are stronger than no explanation, weaker than independently audited
coefficient optima.
\item[Heuristics still lacking an established optimum.] Rotation margin
1.35; 1.25/1.4-pitch run guards; 6\% peaks; $0.6P$ suppression; seam
weights 4, 0.012, 0.02 and closing $0.07P$; 15\% ownership; darkness
0.55; 64 by 640 canvas and network widths/depths; 15\% split and 20\%
sampling; dropout 0.2; body threshold 45\%; cut weight 1.15, clip 3 and
minimum $0.08h_x$; 75-degree continuation; 0.6-pixel simplification;
2.1/2.9-width caps; 12-grid descriptor; drift 0.85/$\pm0.25$; clearance
0.05; maximum join 0.80 x-heights; physical join 0.6 mm; anchoring endpoints
0.24/0.32; search/repair budgets; pen overrun, debounce and timeout settings.
Their operational roles are explained in the relevant chapters. No numerical
optimum is fabricated in place of the other agent's sweeps.
\item[Unresolved physical/source facts.] Actual curve underlying torque
interpolation and assumed friction; exact installed driver sense resistors
and reference voltage; GPIO current-limit source; full rail/ground map;
base pull-down placement and transistor saturation under real load. Calculation
can expose these gaps but cannot replace missing measurements.
\end{description}

\section{Source audit and how to trace assertions}
\label{sec:source-audit}
The named final report, local function bodies, cited original papers and
manufacturer/framework documentation serve different roles. The report is
the source for historical numbers; code is the source for present formulas
and defaults; papers establish method origins; component datasheets establish
conditional electrical relations. A generic citation cannot authenticate an
installed board or as-run checkpoint.

New bibliography entries provide links for the A4988 and carrier current
relation, low-side transistor reasoning, GPIO API, PyTorch conventions,
Wilson/NIST, CTC exposition, RMSprop and source snapshots. Original entries
for Otsu, summed-area tables, thinning, seam methods, LSTM, CRNN, optimisers,
initialisation and edit distance are retained. Sources are cited near the
definitions they support. The substantial worked arithmetic and sensitivity
reasoning in these additions are explicit derivations, not quotations.

Automated full-text access was blocked for some publisher DOI pages and
the ODROID/onsemi pages. Their bibliographic links are retained with access
limitations, rather than claiming they were read successfully. Some older
references have no persistent URL in the original bibliography. This manual
does not certify all forty-four report references or every claim in the
unchanged original 335 pages. In particular the torque paragraph and GPIO
current claim still contain missing-source issues. A later complete reference
audit would need each cited edition/page and the specific claim it supports.
This is deliberately distinct from asserting the report has been repaired.

For a future numerical claim, record input identity, source revision,
checkpoint identity, host, preparation, dtype, units, partition, count,
aggregation and raw evidence location. A bare rounded number cannot supply
those facts. The present expansion adds understanding and traceability while
leaving measurements, source code and the final report unchanged.

\section{Additional vocabulary needed to defend the calculations}
\begin{description}[style=nextline]
\item[Dimensionless coefficient.] A pure ratio or multiplier without units.
Its numerical effect still depends on the scale of the quantity it multiplies.
Changing cost normalisation can change a coefficient's meaning.
\item[Heuristic versus hyperparameter.] A heuristic is a practical rule without
a general optimality guarantee. A hyperparameter is a setting chosen outside
ordinary weight learning. A setting can be both. A learned convolution weight
is a model parameter; dropout probability is a hyperparameter.
\item[Guard, floor, cap and clamp.] A guard handles an invalid or exceptional
case; a floor enforces a minimum; a cap enforces a maximum; clipping enforces
both. A numerical guard prevents arithmetic failure, not necessarily bad input.
\item[Quantile and percentile.] A quantile $q_f$ marks fraction $f$ of a
distribution; a percentile expresses $100f$. The 60th percentile is not
60\% of the maximum. Interpolation conventions can differ on small samples.
\item[Argmax and argmin.] They return the location/index attaining an extreme,
not the extreme value. Ties require a rule, often first occurrence.
\item[Indicator.] $\mathbf1[E]$ is one if statement $E$ is true and zero
otherwise. It encodes a decision without implying a probability.
\item[Recurrence, state and backtracking.] A recurrence computes new entries
from previous ones. A state summarises what a prefix needs for future choices.
Backtracking follows saved choices to recover a path. A recurrence can sum
paths, as CTC does, or optimise one path, as a seam/alignment does.
\item[Dynamic programming versus greedy choice.] Dynamic programming compares
prefix solutions for each state; a greedy rule takes an immediate best choice.
Both can be fast, but they solve different optimisation problems.
\item[Prior and posterior.] A prior expresses information before a particular
observation; a posterior conditions on it. The project's glyph ``prior'' is
a pooled geometric reference, not a fully specified Bayesian probability
distribution. CTC state occupancy is a conditional path posterior under the
model and target.
\item[Likelihood versus confidence.] Likelihood evaluates observations/labels
under a model. A softmax output is normalised model confidence; its match to
real frequencies requires calibration evidence. High confidence can be wrong.
\item[Calibration, two meanings.] Mechanical calibration relates pulse counts
to reference positions/lengths. Probability calibration relates predicted
confidence to observed frequency. The procedures and units are unrelated.
\item[Bias and variance, two contexts.] Measurement bias is systematic offset;
sampling variance is variation across repeated estimates. A neural bias is
an additive learned parameter. An unbiased variance estimator is a statistical
formula, not a guarantee a classifier is fair or accurate.
\item[Regularisation and weight decay.] Regularisation constrains learning
to reduce overfitting. Weight decay shrinks parameters. Dropout, augmentation,
early selection and penalties influence generalisation by different mechanisms;
their strengths cannot simply be added as one scalar percentage.
\item[Epoch, update and iteration.] An update changes parameters after a batch;
an iteration may mean one such processing step; an epoch follows the sampler's
definition of a pass. Replacement/mixture sampling can revisit images within
one epoch. A checkpoint's epoch index is not elapsed time.
\item[Logit and log probability.] A logit is an unrestricted score. Log-softmax
subtracts log-sum-exp to produce log probabilities. Log probabilities are
nonpositive; logits need not be. Exponentiating an arbitrary logit alone is
not a normalised probability.
\item[Jacobian, gradient and chain rule.] A Jacobian contains derivatives of
vector outputs with respect to vector inputs. A scalar loss gradient is the
corresponding vector derivative. Backpropagation efficiently composes these
derivatives without storing every full Jacobian.
\item[Subgradient.] A selected derivative-like value at a nonsmooth point.
ReLU at zero and max ties require conventions; disagreement there is not
necessarily evidence of an incorrect smooth derivative elsewhere.
\item[Float32, float64 and tolerance.] Floating-point types have finite
precision, not exact real arithmetic. Comparison tolerance defines an allowed
discrepancy for a particular scale. A tolerance is neither a p-value nor a
physical instrument uncertainty.
\item[Isotropic versus anisotropic.] Isotropic scaling uses one factor in
both axes; anisotropic scaling uses different factors and changes aspect ratio.
Normalisation describes a transformation, not proof that remaining differences
are exclusively writer style.
\item[Aliasing and interpolation.] Aliasing occurs when sampled detail becomes
indistinguishable from another pattern. Interpolation estimates intermediate
values but cannot reliably recreate information lost by undersampling.
Antialiasing before reduction limits high-frequency content to reduce this risk.
\item[Topology versus geometry.] Topology concerns connections, components and
holes; geometry concerns positions, lengths and angles. Small geometric error
does not necessarily preserve a tiny loop, and a connected skeleton can still
have poor geometric fidelity.
\item[Open-loop versus feedback.] Open-loop axis control estimates position
from commands. Feedback observes a quantity, such as an endstop or pen-up
switch. Having some switches does not make every interior carriage position
closed-loop or verify pen pressure.
\item[Resolution, repeatability, accuracy and uncertainty.] Resolution is the
smallest represented/readable increment. Repeatability describes consistency
under repeated conditions. Accuracy concerns agreement with a reference.
Uncertainty characterises the range attributable to measurement knowledge.
A fine step lattice alone establishes none of the other three.
\item[Macro and micro aggregation.] Macro gives each group equal weight;
micro pools elementary counts. Equal-writer ratios, word means and pooled
character counts therefore answer different questions even on the same files.
\item[Ablation and sensitivity.] Ablation removes a mechanism to assess its
contribution under controlled conditions. Sensitivity varies a setting and
observes response. A sensitivity plot can reveal fragility without proving a
global optimum. No such new execution is part of this edition.
\item[Confusion matrix.] Rows/columns count actual/predicted class pairs.
Row normalisation divides by actual-class counts, showing conditional prediction
fractions. It does not equal class precision, which uses predicted-class totals.
An overall average without counts can conceal class imbalance.
\item[Selection leakage and provenance.] Selection leakage occurs when a test
set influences tuning and is later presented as untouched. Provenance records
where an input, model or number came from. It prevents an old result from being
silently attributed to a newer code revision or different host.
\end{description}

\section{Persistent source links for the retained method references}
These links supplement the original bibliography. Publisher access restrictions
are not treated as successful full-text review. Method citations establish
origins; project arithmetic examples and parameter trade-offs remain explicit
derivations and engineering reasoning.
\begin{description}[style=nextline]
\item[Otsu; integral images.]
\url{https://doi.org/10.1109/TSMC.1979.4310076}\par
\url{https://doi.org/10.1145/800031.808600}.
\item[Thinning; polyline simplification.]
\url{https://doi.org/10.1145/357994.358023}\par
\url{https://doi.org/10.3138/FM57-6770-U75U-7727}.
\item[CTC; CRNN; batch normalisation.]
\url{https://www.cs.toronto.edu/~graves/icml_2006.pdf}\par
\url{https://arxiv.org/abs/1507.05717}\par
\url{https://arxiv.org/abs/1502.03167}.
\item[LSTM original and derivative context.]
\url{https://doi.org/10.1162/neco.1997.9.8.1735}\par
\url{https://www.cs.toronto.edu/~graves/preprint.pdf}.
\item[Dropout; Adam; AdamW.]
\url{https://www.jmlr.org/papers/v15/srivastava14a.html}\par
\url{https://arxiv.org/abs/1412.6980}\par
\url{https://arxiv.org/abs/1711.05101}.
\item[Initialisation and clipping.]
\url{https://proceedings.mlr.press/v9/glorot10a.html}\par
\url{https://arxiv.org/abs/1502.01852}\par
\url{https://proceedings.mlr.press/v28/pascanu13.html}.
\item[Bresenham; Wilson.]
\url{https://doi.org/10.1147/sj.41.0025}\par
\url{https://www.jstor.org/stable/2276774}.
\item[Seam methods and document threshold alternative.]
\url{https://doi.org/10.1109/ICFHR.2014.127}\par
\url{https://doi.org/10.1145/1276377.1276390}\par
\url{https://doi.org/10.1016/S0031-3203(99)00055-2}.
\end{description}
